\chapter*{Conclusion}
\addcontentsline{toc}{chapter}{Conclusion}

This thesis set out to provide a comprehensive, step-by-step framework for modeling financial returns, beginning with univariate modeling techniques and culminating in a multivariate risk-modeling approach based on copulas. Each chapter built on the previous one, combining theoretical underpinnings with practical implementation for key models in financial econometrics.

\hyperref[chap:1]{Chapter 1} introduced the empirical features of asset returns and framed the modeling challenge. It also described the dataset used throughout: daily returns for the S\&P~500 and a 10-year U.S.\ Treasury note index (Bloomberg ticker: SPBDU1BT Index), with a training sample from January~2015 to December~2024 and a forecasting (test) sample from January~2025 to April~2025.

\hyperref[chap:2]{Chapter 2} addressed the first challenge—modeling dynamic volatility—via GARCH(1,1) and GJR–GARCH(1,1). \hyperref[chap:3]{Chapter 3} turned to the parametric modeling of shocks, considering an asymmetric Student’s $t$-distribution and a peak-over-threshold approach from extreme value theory (EVT).

\hyperref[chap:4]{Chapter 4} evaluated model validity and usefulness for risk management. The chosen GARCH specifications captured the conditional volatility structure, and density-forecast diagnostics suggested statistical adequacy in samples. For risk forecasts, however, models based on the asymmetric $t$-distribution failed to achieve acceptable Value-at-Risk (VaR) coverage and independence of violations, while EVT-based models achieved more acceptable coverage but did not pass independence tests.

\hyperref[chap:5]{Chapter 5} proposed a multivariate framework based on a symmetric $t$-copula to capture threshold correlations while, for tractability, abstracting from time variation in dependence. In the test set, this static copula struggled with sudden increases in volatility and, consequently, with independence of VaR violations. Nonetheless, PIT-based diagnostics—histograms and Kolmogorov–Smirnov tests on the normal transform—indicated that the copula-based portfolio density forecasts were statistically well calibrated, supporting their validity as density forecasts. Together with the VaR shortcomings noted above, these findings motivate the broader appraisal of the framework's limitations that follows.

Several limitations should temper the conclusions drawn from this work. First, the empirical scope is deliberately narrow: a bivariate portfolio of the S\&P~500 and a 10-year U.S.\ Treasury index, evaluated on a four-month out-of-sample window (January--April 2025) that, while informative as a stress episode, yields too few violations for high-power backtests and spans only one regime shift. Second, all models are fit once on the full 2014--2024 sample rather than re-estimated through rolling or expanding windows, so the parameters reflect a predominantly low-rate, equity-bull period and may adapt sluggishly to structural breaks. Third, the GARCH(1,1) and GJR--GARCH(1,1) specifications inherit well-known caveats---fixed lag order, near-integrated persistence that complicates long-horizon variance forecasts, and reliance on squared returns as a noisy proxy for conditional variance---while the EVT and asymmetric-t shock models assume time-invariant tail parameters. Finally, the risk analysis is conducted on a hypothetical, equally weighted, buy-and-hold position, abstracting from transaction costs, rebalancing, and liquidity effects that any practical implementation would face. These limitations do not invalidate the framework, but they delimit the conditions under which its forecasts can be trusted; addressing them points to hybrid specifications that pair the conservative tail treatment of EVT with shock processes that admit time-varying skewness and kurtosis, and to asymmetric or dynamic copula families for the multivariate layer.

Throughout the thesis, emphasis was placed on balancing theoretical
clarity with empirical rigor. Modeling choices were always accompanied
by justification, comparisons with alternatives, and diagnostic tools
to evaluate performance. The methodologies were implemented and illustrated using real financial data, with all \Rlogo\ code made openly available in the accompanying \href{https://github.com/ferpmillan/thesis}{GitHub repository} to promote full
reproducibility.

Ultimately, the structured modeling path presented here not only equips
readers with robust and practical tools to approach financial return
modeling, but also underscores the ongoing relevance and vitality of
questions first posed decades ago—questions that continue to drive
innovation and deepen our understanding of financial markets.